\chapter{Hip Pain in Children}

\section*{Definition}
Hip pain in a child may be felt in the groin, thigh, or knee. Causes range from transient synovitis to infection, fracture, inflammatory disease, and slipped capital femoral epiphysis.

\section*{When to See a Doctor}
Urgent assessment is needed for fever, inability to bear weight, severe pain, a hot swollen joint, recent trauma, or a child who appears unwell. Arrange prompt evaluation for persistent, recurrent, unexplained, or function-limiting symptoms.

\section*{How It Develops}
The location of pain may not identify its source: hip disease can be felt in the groin, thigh, or knee. Age, fever, ability to bear weight, range of motion, and the pattern of onset help distinguish transient synovitis, infection, trauma, Perthes disease, and slipped capital femoral epiphysis.

\section*{Causes}
\begin{itemize}
  \item Transient synovitis, septic arthritis, trauma, juvenile idiopathic arthritis, Perthes disease, slipped capital femoral epiphysis, and referred pain.
  \item Medication effects, environmental exposures, and underlying chronic disease may also contribute.
  \item Serious causes are less common but must be considered when symptoms are sudden, progressive, severe, or associated with systemic illness.
\end{itemize}

\section*{Related Symptoms}
\begin{itemize}
  \item Pain, fever, fatigue, nausea, weakness, or changes in normal function
  \item Bleeding, swelling, discharge, breathing difficulty, or neurologic symptoms when a serious cause is present
\end{itemize}

\section*{Medical Evaluation}
\begin{itemize}
  \item History addresses onset, fever, trauma, ability to bear weight, recent infection, and pain location.
  \item Examination includes gait, hip range of motion, the knee and spine, and vital signs.
  \item Testing may include pelvic or hip radiographs, ultrasound, blood tests, or joint aspiration when infection or another serious cause is suspected.
\end{itemize}

\section*{Treatment Options}
Treatment depends on the cause, ranging from observation and analgesia for selected transient synovitis to urgent antibiotics and drainage for septic arthritis or orthopedic treatment for Perthes disease, fracture, or slipped capital femoral epiphysis.

\section*{Self-Care and Home Management}
Avoid known triggers, protect the affected area, maintain appropriate hydration, and record timing and associated symptoms. Use only age- or pregnancy-appropriate medicines recommended by a clinician. Do not delay emergency care while trying home treatment.

\section*{References}
\begin{thebibliography}{99}
\bibitem{hip_pain_in_children:limping} Bono KT, Squires J, Della Morte M. Evaluation of the limping child. \textit{Am Fam Physician}. 2015;91:611--618.
\bibitem{hip_pain_in_children:kocher} Kocher MS, Zurakowski D, Kasser JR. Differentiating between septic arthritis and transient synovitis of the hip in children: an evidence-based clinical prediction algorithm. \textit{J Bone Joint Surg Am}. 1999;81:1662--1670.
\end{thebibliography}
